\subsection{问题二：基于实时排队的剩余路线重规划}

问题二继承问题一路线，锁定已完成项目。实时排队更新模型为
\begin{equation}
q^{real}_{d,i}(t)=\max\{0,q^0_{d,i}(t)+\rho(t)\Delta_{d,i}\},
\end{equation}
其中 \(\Delta_{d,i}\) 来自 APP 实时排队扰动，\(\rho(t)\) 表示扰动随时间回归基准的衰减。

\begin{figure}[H]
\centering
\includegraphics[width=0.88\linewidth]{fig_q2_model_schematic.png}
\caption{问题二锁定已完成项目后的实时重规划流程。}
\end{figure}

\begin{figure}[H]
\centering
\includegraphics[width=0.9\linewidth]{fig_q2_result.png}
\caption{问题二实时重规划的等待节省和得分增益。}
\end{figure}

\begin{figure}[H]
\centering
\includegraphics[width=0.9\linewidth]{fig_q2_validation.png}
\caption{问题二保持原路线与重规划路线对照。}
\end{figure}

实时重规划平均减少等待 \(19.67\) 分钟，平均得分增益 \(3.626\)。共有 \(6\) 个场景改变了剩余项目集合，最大等待节省为 \(119\) 分钟。
